\odiseachapter{Duodécima parte}{La matanza de los pretendientes}

La palabra es \gr{μνηστηροφονία} (mnēstērophonía), derivada de \gr{μνηστήρ} (pretendiente) y \gr{φονία} (matanza, homicidio). El griego solo necesita un vocablo donde el castellano precisa de cinco.

Un solo vocablo para nombrar el asesinato a sangre fría de ciento ocho jóvenes, seguido por una implacable ejecución.

Nolan y todos los directores de cine anteriores a él se las han tenido que ver con este dato. Cierto, los pretendientes eran unos gorrones, y no es menos cierto que algunos de ellos eran enemigos declarados del héroe. Antínoo es un miserable que intenta emboscar a Telémaco, involucrando a otros veinte secuaces. Eurímaco y Ctesipo pagan cara su arrogancia. Agelao intenta organizar a los aterrorizados galanes y se comprende que Odiseo y su gente los eliminen. Pero todos sumados no llegan a treinta, lo que quiere decir que los otros ochenta hombres son masacrados como ganado, sin necesidad y sin piedad alguna. Y el acto que sigue merece mención de honor en el libro universal de la infamia.

No hay forma de que Hollywood muestre semejante salvajada. Pero tampoco era necesario endulzar la historia hasta hacerla irreconocible, empezando por el caso de Telémaco. En su Nodisea, el versionante director insiste en tratar al hijo de Ulises, que se parece a su padre en muchos aspectos ---incluyendo la crueldad---, como un muchachito atribulado e indeciso, además de físicamente torpe.

El resultado es ridículo y arruina una película que cuenta, como hemos visto, con hallazgos excelentes. Cuando los pretendientes intentan emboscar a Telémaco en el templo de Pilos (una escena inventada que no aparece en la épica de Homero, donde el príncipe, avisado por la metomentodo Atenea, simplemente da esquinazo a sus enemigos), tiene que ser papá quien intervenga para proteger al chaval. Peor aún, durante la matanza, papá le ordena que se esté quietecito, no se vaya a manchar las manos de sangre ---da risa imaginarse a Odiseo, rodeado de cien enemigos, preocupado por el bienestar emocional de su mocito--- y a Chris Nodoyuna no le queda otra que mandarlo al piso de arriba, a darse bofetadas con el traidor cabrero Melantio, con el único ánimo de tenerlo ocupado.

No menos caricaturesca es la inverosímil batalla en la que Matt Damon acaba hecho un San Sebastián, a pesar de lo cual Nolan no consigue capturar ni la violencia de la masacre, ni el salvajismo del héroe. Es más, cuando parte de los pretendientes «hinca la rodilla» y se rinde, Ulises les perdona generosamente y con ello nos dan gato por liebre, o, más precisamente, gato por tigre. El Odiseo que nos muestra el canto XXII es exactamente eso: un despiadado tigre de Bengala, decidido a devorar hasta el último de sus enemigos. Convertirlo en un gato de angora es, en mi opinión, un pecado capital y si Nolan no hubiera rodado la escena de las Sirenas, le habría costado penar para siempre en el Tártaro.

%Lo que me lleva al primer capítulo de esta serie. El cine teme a Homero y le teme por culpa de la combinación tóxica de falsa superioridad moral, hipocresía y papanatismo de los clientes de Hollywood, es decir, las modernas sociedades occidentales. Ulises es un personaje difícil de digerir. Valiente, embustero, astuto, implacable, artero, cruel, vil a veces, generoso otras ---no demasiadas---, Odiseo no da muchos agarraderos para dibujarlo como un héroe con el que pueda identificarse un público que el lunes va a ver la \emph{Odisea} y el miércoles come palomitas admirando las aventuras de un chico en pijama llamado Spiderman.
%
%Odiseo es, de hecho, un héroe telúrico, más fácil de identificar con una fuerza de la naturaleza que con un personaje de las operetas Marvel. Un ciclón no es bueno ni malo, una tormenta no tiene principios morales, un volcán en erupción no puede evitar calcinar a todo el que está cerca. Odiseo es todo eso y más y menos. Es también Nadie, el hombre sin nombre, sin familia, ni amigos ni patria ---otra omisión digna del Averno por parte del adormilado director---, el eterno vagabundo, el judío errante, Juan sin Tierra. Es el marino que reproducen Simbad y Nemo, es el amante de la hechicera, es un asesino que ha mandado al Hades a legiones de inocentes y el único que se ha atrevido a escuchar a las sirenas. Es posible inventar mil hombres distintos que se parecen a Odiseo, pero, ¿hacía falta reducirlo a un predicador ambulante?

Pero dejemos que sea Homero quien nos cuente la historia antes de enfrascarnos en el delicado asunto de cómo Hollywood adultera el clásico. En el canto XXI, Penélope invita a los pretendientes a la prueba del célebre arco:

\begin{verse}
«Oídme, altivos galanos, que habéis tomado este techo\\
y un día y otro sin fin estáis comiendo y bebiendo\\
de un hombre que hace ya mucho que falta de él; y pretexto\\
otro no habéis encontrado a decir de palabra que esto:\\
que ansia tenéis por llevarme de esposa y en casamiento.\\
Pues, venga, mis pretendientes, aquí a la vista está el premio.\\
Os dejaré este gran arco del biendivino Odiseo;\\
con el que tense en sus manos el arco en menos esfuerzo\\
y por las hachas, las doce, la flecha la pase en un vuelo,\\
con él me iré aquí dejando mi casa de casamiento,\\
este casal tan hermoso y tan de vivires lleno\\
del que me voy a acordar me creo que hasta en los sueños».
\end{verse}

Toda la escena está muy bien conseguida en el film de Nolan, aunque se salta un detalle esencial.  En la \emph{Odisea}, Telémaco es el primero en intentar tender el arco:

\begin{verse}
[...] y soltó de sus hombros el manto purpurigrana\\
brincando en pie, y de sus hombros soltó la espada aguzada.\\
Hincó primero las hachas abriéndoles antes por cama\\
un largo surco, y dejábalas todas a cuerda alineadas\\
apelmazando la tierra en redor; con asombro miraban\\
lo bien a ley que ordenó ---y antes nunca él lo viera--- las hachas.\\
Y fue a probarse en el arco tras apostarse a la entrada.\\
Tres veces lo hizo temblar, de llevarlo a la gaza en gran ansia,\\
tres de su fuerza falló, esperando las tres en su alma\\
tensar el nervio y pasar por el hierro la flecha lanzada.\\
Y bien lo habría encordado tirando en más fuerza a la cuarta,\\
mas Odiseo de seña negó y reprimiose él la gana.
\end{verse}

Observa, asombrada lectora, perplejo lector, lo que nos dice el Poeta. \emph{Telémaco es capaz de tensar el arco de Odiseo} ---al cuarto intento---, es decir, \emph{es casi tan fuerte como su padre} y, de hecho, casi tan astuto como este. Reprime sus ganas de hacerlo y dice así:

\begin{verse}
«Áyayme, ya un parapoco, un indigno seré, bien lo creo,\\
o a lo mejor soy muy joven y manos aún no las tengo\\
para guardarme de quien me llegó en malhechuras primero.\\
Venid vosotros, que en fuerzas muy por delante yo os veo,\\
haced del arco el intento, y que se cumpla ya el juego».
\end{verse}

El primero de los pretendientes que prueba es Leodes, que, entre los pretendientes, oficia como \gr{θυοσκόος} (thyoskóos), un término que combina \gr{θυο} «ofrenda, incienso» y \gr{σκόος} «observar». Es decir, Leodes es «el que observa las ofrendas» o el adivino del grupo. Además es el más decente de la banda de gorrones:

\begin{verse}
Leodes, hijo de Énope en pie se puso el primero;\\
él les hacía de arúspice y junto a la crátera asiento\\
tomaba siempre al rincón; y era el único allí desafecto\\
a sus excesos, y andaba indignado con todo el cortejo;\\
así que el arco y la viva saeta agarró él el primero.\\
Plantado ya en el umbral, al arco echábale tiento:\\
no lo tensó, que agotó antes sus manos en el tiramiento\\
---sus manos suaves, no usadas---, y dijo a sus compañeros:\\
«Amigos, no puedo armarlo, así que lo tome ya otro».
\end{verse}

Ya hemos dicho que Leodes tiene dotes de adivino, así que más les habría valido a los pretendientes hacerle caso cuando dice:

\begin{verse}
«Amigos, no puedo armarlo, así que lo tome ya otro.\\
A muchos de los mejores este arco les va a dejar rotos\\
el corazón y el aliento, pues es preferible del todo\\
morir a truncos vivir por no ver consumado ese logro\\
que aquí esperando y en corro nos tuvo un día tras otro.\\
Habrá quien guarde en su pecho aún esperanza, afanoso\\
por de Penélope ser, la mujer de Odiseo, el esposo.\\
Mas cuando al arco se atreva y ya lo tenga a sus ojos,\\
mejor que a alguna otra aquea veligalana obsequioso\\
por cortejarla se vuelva; y deje que al matrimonio\\
llame la reina al que, a más de fatal, venga más dadivoso».
\end{verse}

Antínoo, después de burlarse de Leodes, ordena a Melantio que prepare un fuego para untar el arco con grasa y así ablandarlo. El cabrero obedece, pero ni aun así consigue ejecutar la proeza ninguno de la alegre banda. Solo quedan por probar Antínoo y Eurímaco. Entre tanto, Odiseo se escapa al patio y alista a sus dos leales, el boyero Filetio y el porquerizo Eumeo, y les da instrucciones para consumar su plan:

\begin{verse}
Ahora hay que entrar, y se hará uno por uno, no todos a un tiempo;\\
Yo iré primero, vosotros después; y a esta seña id atentos:\\
ninguno de esos galanos de los del ilustre cortejo\\
permitirá que la aljaba y el arco me den para el tiento;\\
pero, por toda la sala el arco llevándolo, Eumeo,\\
me lo pondrás en las manos, diciendo a las siervas tú luego\\
que echen cerrojo a las bien firmes puertas del gran aposento,\\
y si una escucha gemido o estruendo de hombres adentro\\
de nuestros muros, no vaya a querer asomarse allí abriendo,\\
no, que a su jera y labores habrá ella de estar en silencio.\\
Y a ti las puertas del patio te encargo, divino Filetio:\\
ciérralas con el cerrojo echándoles nudo de recio.
\end{verse}

Entre tanto, Eurímaco fracasa con el arco y se queja amargamente.

\begin{verse}
«Áyayme, a fe es por mí mismo y por todos la pena que siento,\\
y, aunque me amarga, no es tanto la boda lo que ahora lamento\\
---hay otras muchas aqueas y sin salir de este suelo en\\
Ítaca milsalinera, y muchas habrá en otros pueblos---,\\
como el que estemos en fuerza quedando todos tan lejos\\
del pardediós Odiseo; ni dar a la cuerda podemos\\
su arco: afrenta que oirán mañana los venideros».
\end{verse}

El detalle no tiene desperdicio. A Eurímaco le preocupa más perder su reputación y que se sepa que ni él ni nadie de la pandilla han podido tensar el arco, que la boda con Penélope. Antínoo, por su parte, ni siquiera lo intenta y echa balones fuera comparando a Odiseo con un dios y, por tanto, dando a su fuerza carácter sobrenatural. Propone, a cambio, seguir con la fiesta. Cuando todos están bastante borrachos ---un detalle importante, ya que es más fácil abatir al enemigo si está ebrio---, Odiseo pone en marcha su treta:

\begin{verse}
«Oíd, los que pretendéis con la bien clara reina alianza,\\
que esto es lo que el corazón desde sus entretelas me manda:\\
a Eurímaco sobre todo y a Antínoo divinaestampa\\
suplico ---y más cuando ha dicho a punto y verdad su palabra---\\
que ahora descansen del arco, que ya los dioses se encargan,\\
y al alba el dios le dará la fuerza a quien se le plazca.\\
Pero dejadme a mí el arco pulido y que prueba yo haga\\
ante vosotros de manos y fuerzas por ver si aún arraiga\\
aquel vigor que en mis miembros retuertos antaño llevaba,\\
o si me lo han ya matado el abandono y la errancia».
\end{verse}

A los pretendientes la propuesta les sienta mal, entre otras cosas por miedo a que un mendigo sea capaz de triunfar donde ellos han fracasado y Antínoo le increpa:

\begin{verse}
Ah, condenado extranjero, en ti no hay ni miaja de seso;\\
¿no te contentas con mesa aquí haber con nosotros, ileso,\\
entre tan alta ralea, sin mengua de yanta, y oyendo\\
nuestra conversa y palabra además? Ningún otro extranjero\\
palabra nuestra llegó aquí a escuchar, ni ningún pordiosero.
\end{verse}

Luego le echa una diatriba, lo llama borracho y le amenaza:

\begin{verse}
Desgracia tanta en ti ahora, si el arco tensas, descuento,\\
pues que de nadie tendrás benevolencia y aprecio\\
en nuestro pueblo, que presto sin más en un negro velero\\
al rey Equeto, de los morideros verdugo y flagelo,\\
te enviaremos: de allí no hay regreso. Así que sereno\\
bebe, y porfía no busques con los que te son más mozuelos.
\end{verse}

A lo que contesta Penélope:

\begin{verse}
«No es noble, Antínoo, ni justo dejar de lo suyo privado\\
a quien llegó a este casal y ganó de Telémaco amparo.\\
¿Crees que si el huésped lograra tensar de Odiseo el gran arco\\
porque en su fuerza fiara y en el vigor de sus manos\\
me iba a llevar a su casa y a hacerme su esposa por pago?\\
A fe que ni él en su pecho está tal cosa esperando.\\
Ninguno se ande por eso hondodolido en su ánimo\\
mientras aquí banquetea, que no, que tampoco es el caso».
\end{verse}

Eurímaco, otro de los pretendientes, mete baza para explicar que ninguno de ellos ---todos hijos de nobles casas--- puede contemplar la posibilidad de que un mendigo arruine su reputación. Penélope le responde, sin pelos en la lengua:

\begin{verse}
«¿Qué buena fama va a haber, Eurímaco, entre la gente\\
para quien casa deshonra comiendo en desprecio los beches\\
de su varón principal? Y por sobra os tomáis que se cuente.\\
El forastero es muy alto, y bien formado y es fuerte,\\
y de hijo ser de buen padre se precia por su progenie.\\
Así que dadle ya el arco pulido por ver si lo tiende.\\
Y cosa más aún os digo que cumpliré si sucede:\\
si lo templara, porque eso a él Apolo por gloria le diere,\\
lo vestiré, con jubón y manteo, de muy buenas vestes,\\
y le daré aguda azcona que a perros y hombres alueñe,\\
y espada doblefileña, y sandalias de buenos cordeles,\\
y su partida armaré adonde el corazón a ir lo lleve».
\end{verse}

No pierdas detalle, biensabida lectora, bienatento lector. A diferencia de lo que nos cuenta Nolan, \emph{Penélope no sabe que el mendigo es su marido}. Y esto es importante porque su intervención favorece los planes de Ulises sin ella saberlo y también lo es porque muestra su generosidad y, sobre todo, \emph{su agencia}. La Penélope de Nolan ayuda a su marido, al que ya ha reconocido, es decir, se comporta como su comparsa. La de Homero tiene voz y voluntad propias, piensa y decide por su cuenta. ¿Es nuestro bardo más feminista que el bienintencionado director?

A Telémaco, en cambio, le falta tiempo para sacar sus galones de hombre de casa:

\begin{verse}
«Ninguno de los aqueos mejor que yo, madre mía,\\
para ofrecer o negar ese arco a quienquier que yo diga,\\
ni de los que señorean por los pedrizales de Ítaca,\\
ni de las islas que miran a Élide belyeguariza;\\
fuerza ninguno me hará malmigrado ni si es gana mía\\
dar a mi huésped el arco y que se lo lleve en propina.\\
Así que ve a tu aposento y de tu labor allí cuida,\\
que es el telar y la rueca, y di a tus mujeres que sigan\\
a su trabajo aplicadas, que el arco cuenta es masculina,\\
muy cuenta mía además, que soy yo quien la casa administra».
\end{verse}

Muy astuto y ya es la segunda vez que nos lo muestra. De un solo golpe consigue mandar a su madre fuera de la estancia ---alejándola así de la inminente carnicería--- y hacerse con el arco, ya sabemos con qué propósito. La pobre Penélope, que no se espera que el hijo querido le meta semejante riña, se marcha apenada y echando de menos al marido ausente, hasta que Atenea, siempre atenta, la deja dormida:

\begin{verse}
Ella marchó a su aposento otra vez, hondamente asombrada,\\
pues las palabras del hijo se entraron sagaces al alma.\\
Ya en las estancias de arriba, adonde subió con sus amas,\\
llora a Odiseo, a su esposo, y lo llora hasta que a ella le lanza\\
bajo los párpados, dulce, el sueño Atenea ojiglauca.
\end{verse}

Eumeo, que como sabemos está en el ajo, se hace con el arma, mientras los pretendientes se burlan de él. Luego le pide a la nodriza Euriclea que eche el cerrojo a las puertas de la sala. Odiseo, entre tanto, empieza a tentar el arco, ganándose también algunas mofas, otra escena muy bien conseguida en la película, lo que demuestra que, si Nolan se queda corto, no es por falta de atención ---ni devoción---, sino por cierta anemia en el coraje.

Y mientras los pretendientes se burlan, Odiseo tensa el arco:

\begin{verse}
Decían los pretendientes; en tanto el ardido Odiseo,\\
bien sopesado y mirado el gran arco por todo su cuerpo,\\
igual que el que hace del canto y de la cítara empleo\\
tensa en la nueva clavija la cuerda muy sin esfuerzo\\
luego de atar una tripa torcida de oveja a su extremo,\\
así, sin pena ni afán, tensó el gran arco Odiseo.\\
Tomó después, por probarlo, entre sus dedos el nervio,\\
y agudo y claro a su mano cantó, como voz de vencejo.
\end{verse}

¡Qué maravilla! Ulises hace cantar el arco «como voz de vencejo» y a Nolan ese detalle no se le escapa, pero se empeña en revestirlo de significado: su héroe hace sonar el arco \emph{siempre} para advertir a sus víctimas de lo que les espera y darles una oportunidad. El del Poeta, por el contrario, carece por completo de escrúpulos y, en gran medida, esa es su ventaja.

Antes de la matanza, el ritual que lo identifica a los ojos de todos:

\begin{verse}
Tomó una vívida flecha que sobre la mesa y en huelgo\\
nuda yacía; guardaba la aljaba en su cuenco agujero\\
a las demás, las que pronto probaran allí los aqueos.\\
Y ya, la vira en su muesca, apuntando tiraba del nervio\\
desde su escaño sentado, y el dardo lanzó, que a derecho\\
pasó por todas las hachas sin yerro alguno en su vuelo\\
de ojo en ojo a través, hasta haber en el patio su término\\
el tiro broncicebado. Y entonces dijo a Telémaco:
\end{verse}

\begin{verse}
«Telémaco, no te avergüenza el huésped que tienes sentado\\
aquí a tu mesa; ni el blanco lo erré ni estirar ese arco\\
largo cansancio me dio: está firme aún el brío en mi brazo,\\
no como, a mi menosprecio, tasábanmelo estos galanos.\\
Hora es de darles ahora también la cena a los dánaos\\
mientras hay luz, que vendrá otro solaz enseguida a animarnos:\\
baile y canción, y el festín en sazón quedará y adornado».
\end{verse}

Humor negro no le falta. Todos sabemos qué cena les espera a los dánaos:

\begin{verse}
Dijo, y de cejas dio seña; y está aguda espada ciñendo\\
ya el de Odiseo del cielo, el hijo querido, Telémaco,\\
y, echando mano a su lanza, junto al escañuelo paterno\\
se planta muy bien armado de un bronce como de fuego.
\end{verse}

Nada de «tú quieto, hijito, que papá se ocupa». Telémaco es un héroe por derecho propio: quiere defender a su familia, vengar las afrentas recibidas y recuperar su hacienda, igual que su padre. No necesita este guerrero, «armado de un bronce como de fuego», que nadie le ofrezca un \emph{safe space}.

Al contrario que en la Nodisea, el original no hace un melodrama de la muerte de Antínoo, ni falta que hace. De hecho, es el primero que cae:

\begin{verse}
Y enderechó, esto diciendo, a Antínoo una amarga saeta.\\
Estaba aquel levantando en eso una copa muy bella\\
de oro y ansas gemelas ---la asían sus manos apenas---\\
para su vino beber, y a matanza su alma era ajena.\\
¿Cómo iba allí uno a pensar que entre aquellos sentado a la mesa\\
un hombre, solo entre muchos ---por más que sobrado de fuerza---,\\
la mala muerte vendría a traerle y el ángela negra?
\end{verse}

Odiseo lo despacha mientras da un trago de vino, sin darle tiempo a rechistar. Como es marca de la casa, la descripción de su muerte es exacta y terrible, mucho más, en mi opinión, que la manida escena alternativa que nos da Nolan:

\begin{verse}
A él encarando, la gorja Odiseo le abrió con la flecha,\\
y por el mórbido cuello salió la punta a derechas.\\
Venciose a un lado, y la copa cayó de sus manos sin fuerzas,\\
y un caño espeso bajaba por su nariz de carrera\\
de humana sangre, y con un golpe brusco del pie echó la mesa\\
lejos de sí, y el manjar que tenía esparció por la tierra:\\
el pan, las carnes doradas en sangre y polvo. Se encrespan\\
los pretendientes, al verlo caer, en barbulla violenta,\\
y saltan de sus sitiales, y en pánico por la tarbea\\
miran, buscando, alredor los firmes muros de esa;\\
mas que tomar no había escudo ninguno ni lanza cereña.\\
Y con palabra enconada así a Odiseo lo increpan:
\end{verse}

\begin{verse}
«A malas has disparado, extranjero; pero otro torneo\\
no te vendrás a encontrar: ya has ganado el barranco funesto.\\
Pues ¿qué?, si has muerto al que era el más preclaro mancebo\\
de Ítaca belmoceril: de los buitres serás alimento».
\end{verse}

Todavía no se han dado cuenta de la que les espera. El torvo héroe los saca de su error enseguida:

\begin{verse}
«Perros, pensabais que nunca de vuelta yo a casa llegara\\
de aquella tierra troyana, y así mi casal devorabais,\\
y por la fuerza os entrabais al lecho con mis esclavas,\\
y a mi mujer en asedio ---y yo vivo--- la cortejabais\\
sin miedo a los que en el cielo morada tienen bien ancha\\
ni a que de hombre viniera andando el tiempo venganza:\\
ahora atado lleváis el nudo de la matanza».
\end{verse}

Eurímaco, entonces, intenta salir del apuro echándole la culpa de todo al caído Antínoo y ofreciendo compensar los desmanes causados:

\begin{verse}
«Si es la verdad, Odiseo de Ítaca, que has regresado,\\
razón te sobra en lo dicho de cuanto hicimos los dánaos,\\
desmanes mil en tu casa, y mil también en el campo.\\
Mas el que culpa ha de todos lo tienes ahí derrumbado,\\
Antínoo; que era este solo el que nos empujaba a esos casos,\\
no porque boda anhelara o necesitárala tanto,\\
sino mirando otro fin que no quiso el Cronión consumarlo:\\
ser en Ítaca él, la de firme solar, soberano\\
luego de armarle a tu hijo celada para matarlo.\\
Él yace ahora en su suerte, perdona tú a tus vasallos;\\
que del común ya te haremos arreglo después en reparo\\
de cuanto se haya bebido y comido aquí en tu palacio,\\
y cada uno un desquite en veinte bueyes tasado\\
te pagaremos en bronce y en oro hasta haber aplacado\\
tu corazón: y hasta aquí no es baldón que tú sigas airado».
\end{verse}

Se trata de una escena que Nolan ---es decir, Hollywood--- nos escamotea. ¿Por qué? Porque lo razonable, \emph{lo honorable}, desde nuestra óptica moderna, sería perdonar a los enemigos. A fin de cuentas, Odiseo se ha vengado del traidor, ha dado un escarmiento y le ofrecen compensación por los desmanes cometidos.

Nada de eso. Al héroe nadie va a privarlo de su venganza:

\begin{verse}
Y a él, mirándolo torvo, le dijo el ardido Odiseo:\\
«Eurímaco, ni aunque pagarais con todos los bienes paternos,\\
cuantos ahora tenéis, y hasta más que añadierais a esos,\\
ni así pondría a mis manos de la matanza yo el freno\\
hasta que la demasía pagado haya bien el cortejo.\\
Pelear de frente es ahora lo que a vosotros os dejo\\
o huir si a muerte y malhado alguno ve escapamiento;\\
mas se me da que ninguno va a hurtarse del trance funesto».
\end{verse}

Eurímaco se da cuenta de que lo único que les queda es pelear o morir y dice:

\begin{verse}
«Amigos, no detendrá este hombre su mano en holganza,\\
que ahora que dueño se ha hecho del arco pulido y la aljaba\\
desde el umbral reluciente nos flechará hasta que haga\\
carnicería de todos: pensemos ya en la batalla.\\
Desenvainad las espadas, poned vuestras mesas de cara\\
a las saetas de pronto morir, y vayamos en masa\\
todos a él, y de umbral y de puertas hagamos que salga\\
para llegarnos por la ciudad y que vuele la alarma;\\
así ya ese hombre hoy por última vez con el arco tirara».
\end{verse}

Atención, los pretendientes no necesitan que el cabrero les tire las armas desde el piso de arriba, como ocurre en la película. Ya están armados, llevan a cuestas sus espadas de bronce y en ese sentido Homero es más atrevido ---por enésima vez--- que Nolan. El tímido director pretende que los dánaos no tienen apenas con qué defenderse, lo que hace la hazaña de Odiseo menos espectacular ---aunque quizás más creíble a ojos de un espectador moderno---. El Poeta, sin embargo, no solo los arma, sino que les da la voluntad de defenderse. Y la verdad, uno se pregunta cómo se las va a arreglar el furioso arquero para despachar a un centenar de jóvenes blandiendo espadas y protegiéndose de sus flechas con improvisados escudos de madera. La lectura, entre líneas, está muy clara. La hazaña solo es posible porque los pretendientes son humanos y Ulises, aquí, tiene algo ---mucho--- de divino. Y no es un dios misericordioso. Odiseo, el de las muchas caras, es también el dios de la muerte y la destrucción.

Eurímaco, de hecho, muere con honor, intentando defenderse:

\begin{verse}
A la que dijo él así, sacó su espada buída\\
a filo y filo aguzada, broncina, y echósele encima\\
a hórrido grito a la vez que el divino Odiseo una vira\\
le disparaba alcanzándole el pecho bajo la tetilla:\\
raudo en el hígado el tiro se hundió; y la espada que asía\\
se le escurrió entre los dedos, y desmadejado él caía\\
enguruñido sobre una trabanca, y rodó la comida\\
y el cáliz doble de cuenco; y en estertor de agonía\\
golpeaba el polvo su frente, y por detrás sacudía\\
con ambos pies el sitial; y la niebla veló sus pupilas.
\end{verse}

El siguiente en abrazar la parca es Anfínomo y no es Odiseo quien lo mata, sino Telémaco:

\begin{verse}
Lanzose entonces Anfínomo hacia el glorioso Odiseo\\
bien a derecho ---en su mano llevaba el alfanje malfiero---\\
por si apartarlo de puertas pudiera. Mas antes Telémaco\\
fue por detrás y, arrojándolo, hincole el venablo bronceño\\
entre hombro y hombro a mitad, y se lo pasó por el pecho:\\
él al caer retumbó, y con toda su frente dio al suelo.
\end{verse}

Inmediatamente, el de Odiseo corre junto a su padre y le dice:

\begin{verse}
«Padre, un escudo y dos lanzas ahora voy a traerte\\
y yelmo todo de bronce bien ajustado a tus sienes,\\
y me armaré yo también, y a porquero y boyero de bueyes\\
daré sus armas: ceñirnos los cuatro ahora bien nos conviene».
\end{verse}

Así que Homero no pretende que sea Ulises quien luche en solitario, pues le acompañan en la refriega su hijo y sus fieles. El héroe, que sabe lo suyo de batallas, le dice:

\begin{verse}
«Tráelas corriendo mientras me defienden las flechas que quedan,\\
no vaya a ser que al dejarme aquí solo me ganen la puerta».
\end{verse}

\begin{verse}
Tal dijo; y obedeciendo el muchacho a su padre, se marcha\\
al almacén al que habían llevado las ínclitas armas.\\
De allí sacó cuatro escudos, y cuatro parejas de lanzas,\\
y cuatro yelmos de bronce con sus cimeras crinadas;\\
con ellos raudo salió, y raudo al padre llegaba.\\
Él mismo el bronce a su cuerpo ciñó muy antes que nada,\\
y, como él, los dos siervos vistieron sus armas galanas,\\
y del ardido Odiseo a redor, mienteviva, se plantan.
\end{verse}

Hay una escena en la secuencia de la matanza, en Nolan, que da casi risa. Matt Damon, al que le falta un poco de envergadura para encarnar al divino Odiseo ---yo le habría sugerido considerar a Henry Cavill, que da más el tipo---, se defiende sin más arma que una flechita, gesticulando mucho para apartar a los enemigos que le acosan, tan pánfilos como él. Homero nos pinta una batalla mucho más fiera:

\begin{verse}
Mientras duráronle a este para defenderse los dardos,\\
de los galanos fue haciendo blanco en la sala tras blanco,\\
y siempre a uno mataba, y caían amontonados.\\
Mas cuando ya al rey flechero las flechas lo abandonaron,\\
del firme muro en la sala contra un pilar apoyado,\\
por la crujía de todo esplendor dejó quieto el arco,\\
luego a los hombros echose escudo cuatriadobado,\\
caló en su testa robusta el yelmo bienaforjado,\\
corcelicrino ---de arriba pendía espantable penacho---,\\
y, rematados en bronce, asió dos recios venablos.
\end{verse}

¿Cómo te sentirías, asustada lectora, petrificado lector, si vieras delante de ti a un gigantón furioso, enfundado en bronce, armado hasta los dientes? Pero en esto, el cabrero Melantio consigue hacerse con no pocas lanzas y escudos, lo que desequilibra todavía más el ya desigual combate:

\begin{verse}
Y el cabrerón de las cabras, Melantio, marchó, esto diciendo,\\
por las lumbreras traseras al almacén de Odiseo.\\
Sacó de allí doce escudos, y mismas lanzas con ellos\\
y mismos fueron y en bronce corcelicrinos los yelmos.\\
Y vuelve a todo correr por repartirlo al cortejo.\\
Le flaquearon las piernas y el corazón a Odiseo\\
a la que vio que, ya de armas ceñidos, estaban blandiendo\\
largo rejón en sus manos: trabajo iba a ser no pequeño.
\end{verse}

Cuando Melantio se dirige al almacén para seguir trayendo armas, Ulises manda a Eumeo y Filetio para que lo intercepten. No les da orden de matarlo, sin embargo, sino de dejarlo atado de pies y manos y colgado del techo ---las razones ya nos las imaginamos---. Así que no es Telémaco, como en la Nodisea, quien enfrenta al cabrero, sino los dos fieles sirvientes, que tras cumplir su misión se unen al héroe y su hijo. La cosa, entre tanto, se ha puesto fea. Son cuatro contra cien.

\begin{verse}
Ellos vistieron sus armas, cerraron la puerta brillante,\\
y hacia el sagaz Odiseo volvieron, el milmañasartes.\\
Estaban cuatro, alentando furor, allí en los umbrales,\\
y enfrente, dentro de casa, turba de muchos pujante.
\end{verse}

Así las cosas, es inevitable que Atenea aparezca:

\begin{verse}
Atena, la hija de Zeus, a aquellos vino a allegárseles\\
a Méntor siendo en figura y a par en voz semejante.\\
Y se alegró de esa vista Odiseo y le dijo al instante:
\end{verse}

\begin{verse}
«Méntor, ayuda a este estrago, recuerda a tu amigo querido,\\
quien tantos bienes te hacía, el par en años contigo».\\
Dijo, aun sabiendo que era Atena la Salvadora.
\end{verse}

Uno de los pretendientes, Agelao, sin saber que tienen delante a una diosa, amenaza al supuesto Méntor ---sin darse cuenta tampoco, aparentemente, de que ha aparecido de la nada---:

\begin{verse}
Enfrente los del cortejo gritaban en una voz sola;\\
y a ella le lanza reproche el primero Agelao Damastórida:\\
«Méntor, que no te engatuse con ese su verbo Odiseo\\
para luchar contra los pretendientes y a él socorrerlo.\\
Que nuestra empresa vendrá a terminar como ahora te cuento:\\
tan pronto como a esos dos ---al padre y al hijo--- matemos,\\
muerto con ellos caerás tú también, que es muy caro el empeño\\
que te propones aquí; tu cabeza será nuestro precio.\\
Y cuando a bronce os hayamos quitado las fuerzas y aliento,\\
cuantas habiendas posees ---las que están bajo techo y a cielo---\\
las juntaremos con las de Odiseo; y ni dejaremos\\
en tu palacio vivir a tus hijos como herederos,\\
ni que tu viuda y tus hijas por Ítaca anden sin dueño».
\end{verse}

Atenea se enfada con estas palabras, pero no con los galanes, que considera carne de cañón, sino con su protegido Odiseo, por pusilánime. Claramente los estándares de los olímpicos no son como los de los humanos:

\begin{verse}
«Ya en ti no están, Odiseo, el brío firme y pujanza\\
de cuando tú por la bien blanca Helena, la bienempadrada,\\
por nueve años luchaste contra los teucros sin tasa\\
y a muchos hombres mataste en la penosa batalla\\
y por tu ardid se tomó la ciudad de su rey calliclara.\\
¿Cómo es que ahora, cuando has alcanzado tu hacienda y tu casa,\\
plantado ante estos galanos, tan fuerte tú, te amilanas?\\
Ah, tente aquí, melón dulce, a mi par y mira mis ganas,\\
que vas a ver cómo son ante cáfila hostil las agallas\\
con las que paga este Méntor Alcímida gracias pasadas».
\end{verse}

Y en otro de sus giros desconcertantes, la diosa, después de semejante bravuconada, se convierte en golondrina y se instala en una viga del techo desde donde seguir cómodamente la refriega. A ojos humanos, simplemente, Méntor desaparece. Pero los aqueos, enzarzados en mortal combate, no están para trucos de magia. Agelao ve su oportunidad de ganar la batalla:

\begin{verse}
Buscó Agelao Damastórida en los pretendientes arrojo,\\
y Anfimedonte también, y Eurínomo, y aun Demoptólemo,\\
y el de Políctor, Pisandro, y el ardentísimo Pólibo,\\
a la sazón los más bravos, muy por encima de todos,\\
de los que, vivos aún, por sus vidas peleaban; que a otros\\
el arco y ristra de dardos rodar los hizo ya al polvo.\\
Y fue Agelao quien tomó la palabra y habló de este modo:
\end{verse}

\begin{verse}
«Amigos, pronto tendrá este hombre su mano en holganza;\\
Méntor ya veis que no está, que dijo sus baladronadas\\
y los dejó en los lumbrales igual de solos que estaban.\\
Así que no tiréis todos la larga lanza a la brava,\\
que tiren seis los primeros, a ver si Zeus nos regala\\
herida hacer a Odiseo y prez ganar de la fama.\\
No hay de los otros cuidado a la que a tierra ese caiga».
\end{verse}

Agelao tiene muy claro que si abaten a Odiseo, la batalla se acaba. Seis de los pretendientes le arrojan entonces las lanzas, pero Atenea las desvía.

\begin{verse}
Dijo, y lanzaron entonces los seis, como había mandado,\\
en ansia viva; e inanes Atena dejó esos venablos.\\
Uno en la jamba no más de la bien firme sala hizo blanco,\\
otro en la puerta de sólido engarce, y de otro fue el dardo\\
broncicebado a golpear en el muro y caer rebotando.
\end{verse}

Hay dos lecturas posibles aquí. Una es admitir que incluso un asesino profesional como Odiseo no podía derrotar a cien jóvenes armados y, en consecuencia, debe su victoria exclusivamente a la intervención divina. Es una interpretación que circula, como hemos visto, en toda la \emph{Odisea} y, en general, en toda la épica griega. Los dioses, en último término, son los que deciden y si Ulises sale victorioso de la refriega no es por valiente, hábil o astuto, sino por ser el protegido de Atenea.

Otra lectura, que yo siempre he preferido, es asumir que Homero usa a menudo a los dioses en el sentido metafórico. Zeus nos da una pista durante el concilio de los olímpicos nada más arrancar la \emph{Odisea}:

\begin{verse}
Ay, los humanos siempre andan echando a los dioses la culpa;\\
que es que les van de nosotros sus males repiten, y nunca\\
ven que no estaba en su sino el revés por sus propias locuras.
\end{verse}

A la inversa, quizás los humanos también «echan la culpa» a los olímpicos cuando la suerte les favorece. Así que quizás no hace falta recurrir a la diosa ojiglauca para frenar las lanzas. Quizás, simplemente, Odiseo es un tipo afortunado, como se dirá a sí mismo, años más tarde, viejo y desengañado en Ítaca:

\begin{verse}
Al menos ahora eres lo bastante lúcido\\
para no engañarte, como solías hacer\\
cuando te imaginabas ser el gran Ulises,\\
el héroe y el viajero y el amante de Circe.\\
Cuando recuerdas, todo lo que ves\\
es ruido y furia. Un matón cruel y astuto\\
que tuvo más suerte que los demás.
\end{verse}

El caso es que las lanzas no aciertan y ahora es el turno de Ulises y los suyos:

\begin{verse}
Cuando los cuatro las lanzas de los del cortejo esquivaron,\\
habló el sufrido Odiseo a los que tenía a su lado:\\
«Diría, amigos, que es vez de que pica arrojemos nosotros\\
sobre el tropel de los pretendientes que están tan ansiosos\\
de coronar anteriores oprobios con nuestro despojo».
\end{verse}

Por supuesto, ni Odiseo ni sus fieles fallan.

\begin{verse}
Dijo; arrojaron entonces sus lanzas vibrantes ya ellos\\
tirando al frente; e hincó a Demoptólemo muerte Odiseo,\\
a Euríades muerte Telémaco, a Élato muerte el porquero,\\
e hincó el boyero de bueyes muerte a Pisandro. Mordieron\\
al mismo tiempo los cuatro allí el sinfínido suelo,\\
y recularon buscando un rincón los demás del cortejo.\\
Y ellos lanzáronse ahí a arrancar el astil de los muertos.
\end{verse}

Los pretendientes tiran de nuevo, y esta vez Telémaco y el porquero reciben heridas leves:

\begin{verse}
Los pretendientes sus lanzas por vez segunda arrojaron\\
en mismas ansias, y Atena frustró casi todos sus dardos.\\
Uno en la jamba otra vez de la bien firme sala hizo blanco,\\
otro en la puerta de sólido engarce, y de otro el venablo\\
dio, ponderoso de bronce, en el muro y cayó rebotando.\\
Anfimedonte, no obstante, a Telémaco hiriolo en la mano,\\
que le rozó la muñeca, y le hizo por cima desgarro.\\
Ctesipo a Eumeo pasole por sobre el escudo su dardo\\
que, haciendo araño en su hombro, siguió hasta el suelo volando.
\end{verse}

Inevitable imaginar ahora a Telémaco llegando a casa con la muñeca vendada y diciéndole a su madre: «deberías haber visto al otro tipo».

\begin{verse}
Y ya otra vez Odiseo sagaz y su corro arrojaron\\
sus azagayas punteras sobre el tropel de galanos.\\
Asolador Odiseo en Euridamante hizo blanco,\\
Telémaco en Anfimedonte, en Pólibo el guardamarranos,\\
y por su parte a Ctesipo el que tiene de bueyes cuidado\\
le hincó la suya en el pecho mientras decíale ufano:
\end{verse}

\begin{verse}
«Politersida, el amigo de befas, no más te abandones\\
a altisonancias de poco saber, es mejor que a los dioses\\
dejes el verbo, que pueden los dioses muy más que los hombres.\\
De aquella pata de vaca adehala es en torna este golpe,\\
la que le diste a Odiseo mendigo en sus propios salones».
\end{verse}

Filetio se refiere a lo que ha ocurrido un rato antes, durante la cena, cuando Ulises aparece disfrazado de mendigo. Ctesipo, hijo de Politerses (de ahí «Politersida»), le arroja una pata de vaca al supuesto pedigüeño, que este esquiva sin esfuerzo. El boyero le devuelve la broma con un lanzazo en el pecho.

Las cosas, a partir de aquí, se ponen feas para los pretendientes, entre los que cunde el pánico, haciéndoles romper filas. La refriega se convierte en carnicería:

\begin{verse}
[...] Al punto Odiseo\\
al de Damástor herida de lanza le abrió en cuerpo a cuerpo.\\
Luego Telémaco a Leócrito hirió, al hijo de Evénor,\\
con su azagaya en mitad de la ijada y de claro entró el rejo:\\
aquel de bruces cayó, y con toda su frente dio al suelo.\\
Entonces ya alzó Atenea desde lo alto del techo\\
su adarga malamuertera, y el pánico se hizo aleteo.\\
E iban de acá para allá como vacas que, al pasto saliendo,\\
zumbón el tábano espanta metiéndose en ellas molesto,\\
por las sazones de abril, cuando largos los días van siendo.\\
Como los buitres zarpirrecorvos y piquirretuertos\\
se echan encima de las avecicas desde los oteros\\
y ellas se encogen bajo las nubes y a ras van de suelo\\
y ellos saltando sobre ellas las matan, pues ni abrigadero\\
para esas hay ni escapada, y el cobro alegra al montero,\\
así, irrumpiendo en la sala, los cuatro a los del cortejo\\
van golpeando a voleo en redor, vuelan gritos horrendos\\
de las cabezas quebradas, y humea de sangre ya el suelo.
\end{verse}

Y el Poeta, por si no lo teníamos claro, subraya la crueldad del héroe:

\begin{verse}
Leodes, corriendo hacia él, las rodillas tomó de Odiseo,\\
y suplicándole así le decía en alígeros verbos:\\
«De hinojos ruego, Odiseo, y tú tenme pena y reparo:\\
a fe que nunca a ninguna mujer de las de tu palacio\\
ofensa dije ni hice, que siempre a los otros galanos,\\
si alguno cosa así hacía, trataba yo de frenarlos.\\
Mas ni con eso logré sujetar la maldad de sus manos,\\
y así, por su altanería, funesto sino alcanzaron.\\
Yo, que un arúspice he sido entre ellos y nada hice malo,\\
¿he de caer? ¿Del bien hecho hasta ahora no hay gracia ni pago?».
\end{verse}

Recuerda, memoriosa lectora, sesudo lector, que Leodes, el de las finas manos, es un buen tipo, pues el propio Poeta nos dice que era el único a quien las iniquidades resultaban odiosas y que se indignaba con todos los pretendientes por gorrones y maleducados. ¿Lo perdonará Ulises?

\begin{verse}
Y lo miró torvifiero el ardido Odiseo y le dijo:\\
«Ya que te jactas de haber sido tú su oficiante adivino,\\
seguro, sí, que mil veces habrás en la sala pedido\\
porque mi dulce regreso lejos quedara, y contigo\\
mi esposa amada se fuera y te alumbrara a ti hijos;\\
por eso no escaparás a este final maldolido».
\end{verse}

\begin{verse}
Tal que hubo dicho, una espada tomó entre sus túrgidos dedos\\
que allí yacía ---la había soltado Agelao por el suelo\\
cuando murió---, y con ella segole a mitad el pescuezo.\\
Y su cabeza rodó por el polvo, aún balbuciendo.
\end{verse}

Nada que ver con el asaeteado San Cristodiseo que nos vende Nolan. Lo que me lleva, en mitad de la refriega, a una reflexión importante. Todo personaje de la literatura puede reinterpretarse si se quiere ---la literatura quizás no sea otra cosa que escribir y reescribir siempre las mismas historias--- pero a un precio. Es posible imaginar ---se ha hecho--- una novela en la que Judas Iscariote es un santo y Jesús de Nazaret un pobre hombre corto de luces, pero hay que saber a lo que se arriesga uno. Cuando uno se empeña en vestir un lobo con piel de oveja, el riesgo es caer en la caricatura.

%Sobre todo porque Homero parece empeñado en ofrecernos una complejidad inabarcable. Otro que se las está viendo negras es Femio, el bardo, del que sabemos que es inocente y no ha cometido otro pecado que el de cantar, a la fuerza, para los galanes.
%
%\begin{verse}
%Trataba aún de esquivar el negro trance el Terpíada,\\
%Femio, que allí por la fuerza cantaba para esa cuadrilla.\\
%Estaba en pie sosteniendo clarisonera su cítara\\
%junto al portillo, y dos cosas en mientes con él revolvía,\\
%si de la sala escapar y allegarse al altar que tenía\\
%Zeus del Hogar bienobrado en el patio ---quemado allí habían\\
%siempre Odiseo y Laertes tasajos mil de vaquillas---,\\
%o si acercarse a Odiseo de hinojos y en rogativa.\\
%Tal cavilando, en la duda, pensó que mejor le sería\\
%tomar allí al de Laertes en ruego por las rodillas.\\
%Así que puso en el suelo entonces la cuéncava cítara,\\
%entre un sillón tachonado de plata y una alta vasija,\\
%y él mismo fue ante Odiseo y se abrazó a sus rodillas,\\
%y, suplicando, palabras aladas a él le decía:
%\end{verse}
%
%\begin{verse}
%«De hinojos ruego, Odiseo, y tú tenme pena y respeto:\\
%tendrás después en ti mismo un pesar si por caso a un aedo\\
%matas, a mí, que a los dioses y hombres en canto celebro.\\
%Maestro propio me soy, y el dios ha insuflado en mi pecho\\
%senda de toda canción, y sé bien que cantarte a ti puedo\\
%como si fueras un dios: no te empeñes en darme degüello.\\
%Y hasta Telémaco, tu hijo querido, diríate esto:\\
%que no venía a tu casa por gusto ni propio deseo\\
%para cantarles a los pretendientes en sus banqueteos;\\
%siendo ellos más, y más fuertes, traíanme aquí no queriendo».
%\end{verse}
%
%La absolución parece justa aquí. Pero el Poeta nos da a entender que Femio solo salva el pellejo gracias a Telémaco, que le dice a su padre:
%
%\begin{verse}
%«Tente, y a bronce no hieras al que es culpable de menos;\\
%y salva, a más, al heraldo Medonte, que siempre con celo\\
%en nuestra casa cuidaba de mí desde que era pequeño;\\
%eso si no lo ha matado Filetio o el mismo porquero,\\
%o se topó por la casa contigo, furioso tú y ciego».
%\end{verse}
%
%«O se topó por la casa contigo, furioso tú y ciego». Odiseo un vendaval, un volcán en erupción, una fiera hambrienta, una máquina de matar, un golem enloquecido al que solo puede controlar ---apenas--- su hijo. Y solo gracias a él Odiseo indulta a Femio y Medonte. Dice el milardido asesino:
%
%\begin{verse}
%«No tengas pena, que ya él te ha guardado y a salvo te ha puesto;\\
%sábelo en tu corazón, y que puedas decirle a otro luego\\
%cuánto a la maleficencia aventaja el obrar en lo bueno.\\
%Ahora salid de la casa y buscaos afuera un asiento\\
%lejos de muerte en el patio, tú y el cantor de mil cuentos,\\
%mientras por estas estancias termino el trabajo que debo».
%\end{verse}

Terminada la matanza, el lobo en cuestión recorre la sala, por si queda alguien por rematar.

\begin{verse}
Miró Odiseo por toda la estancia por si un hombre vivo\\
hallaba allí agazapado esquivando su muy negro sino.\\
Mas vio que ya estaban todos en sangre y polvo caídos.\\
Muchos. Que igual que esos peces que dan a cala en abrigo\\
los pescadores sacándolos fuera del mar espumizo\\
a la su red enmallados, y sobre la arena esparcidos\\
quedan, echando de menos las olas del mar, reunidos\\
mientras les quita la vida el haz del sol con su brillo,\\
así yacía el tropel de galanos, en parva abatidos.
\end{verse}

Cuando la anciana entra en la habitación, el Poeta nos muestra en qué se ha convertido su bebé:

\begin{verse}
Halló a Odiseo en mitad de los cuerpos que ha nada abatieran\\
embadurnado de sangre y matanza cual si un león fuera\\
que de comer hasta hartarse de un toro campero volviera:\\
el pecho todo ---y a un lado y a otro sus dos carrilleras---\\
rojo de sangre, espanto de ver al mirar de cualquiera;\\
tal, pringoriento Odiseo por brazo arriba y por piernas.
\end{verse}

Imagina, asustada lectora, trémulo lector, que eres tú quien contempla al criminal embadurnado de sangre. Posiblemente saldrías corriendo, caerías de hinojos entre arcadas o te desmayarías. Euriclea, por el contrario, se pone a gritar de alegría.

\begin{verse}
Apenas vio los cadáveres ella y la sangre sin cuenta,\\
rompió a gritar de alegría, que vio muy grande proeza,\\
pero Odiseo del ansia la tuvo y su gana refrena,\\
y a ella, alzando la voz, palabra en alas le lleva:\\
Por dentro alégrate, vieja, y el grito contén de alegría;\\
gloriarse sobre los muertos ni es santo ni va a ley divina.
\end{verse}

Parece que la masacre ha terminado, ¿no? Pues aún falta lo peor. Pero de eso nos ocuparemos en el siguiente capítulo.
